\begin{algorithm}[htbp]
\caption{Skip Module Hyperparameter Selection}
\label{alg:hyperparam}
\begin{algorithmic}[1]

\Require Parameter space $\Theta$, validation set $D_{\text{val}}$,
DL model $\mathcal{M}$, recall tolerance $\delta_R$,
overhead ratio $\beta$, weights $w_S, w_R$
\Ensure Optimal parameters $\theta^*$

\State baseline $\gets$ \textsc{EvaluatePipeline}($D_{\text{val}}$, skip=None, $\mathcal{M}$)
\State Extract $R_{\text{base}}$, $T_{\text{DL}}$ from baseline
\State best\_score $\gets -\infty$; \quad $\theta^* \gets$ None

\For{each $\theta \in \Theta$}
\State results $\gets$ \textsc{EvaluatePipeline}($D_{\text{val}}$, $\mathcal{S}(\theta)$, $\mathcal{M}$)
\State Extract $R(\theta)$, $S_r(\theta)$, $T_{\text{skip}}(\theta)$ from results
\If{$R(\theta) \geq R_{\text{base}} - \delta_R$
\textbf{ and }
$T_{\text{skip}}(\theta) \leq \beta \cdot T_{\text{DL}}$}
\State $\rho(\theta) \gets 1 - \dfrac{R_{\text{base}} - R(\theta)}{\delta_R}$
\State score $\gets w_S\, S_r(\theta) + w_R\, \rho(\theta)$
\If{score $>$ best\_score}
\State best\_score $\gets$ score
\State $\theta^* \gets \theta$
\EndIf
\EndIf
\EndFor

\State \Return $\theta^*$

\end{algorithmic}
\end{algorithm}

The skip module $\mathcal{S}$ contains several hyperparameters (e.g., motion
threshold, accumulation window size) that can significantly impact the
performance of the overall system. To identify the optimal configuration, we
employ a grid search based hyperparameter optimization procedure on a validation
set $D_{\text{val}}$ derived from our video dataset with respect to several
system-wide metrics defined in Section~\ref{sec:metrics}. More specifically, to
select the optimal hyperparameters for the proposed skip module, we employ a
constrained optimization procedure on the validation set $D_{\text{val}}$. Let
$R_{\text{base}}$ denote the baseline recall, i.e., the fraction of fire/smoke
frames correctly detected by the pipeline without skipping, and let $T_{\text{DL}}$
denote its mean per-frame inference time. For each candidate parameter set
$\theta \in \Theta$ obtained by grid search, we evaluate the full pipeline
$\text{Read} \rightarrow \mathcal{S}(\theta) \rightarrow [\mathcal{M}]$ and
compute three metrics: recall $R(\theta)$ (the fraction of fire/smoke frames
correctly detected), skip rate $S_r(\theta)$ (the fraction of ground-truth
negative frames skipped by $\mathcal{S}$), and mean per-frame skip module time
$T_{\text{skip}}(\theta)$. Formal metric definitions are given in
Section~\ref{sec:metrics}.

The skip rate $S_r(\theta)$ is computed on $D_{\text{val}}$ following the
definition in Section~\ref{sec:metrics}, with $N_{\text{neg}}$ instantiated as
$N_{\text{neg}}^{\text{val}} = |\{f_i \in D_{\text{val}} : y_i = 0\}|$, which
is fixed by the validation set and independent of $\theta$.

The mean per-frame inference times are defined as:
\begin{equation}
    T_{\text{DL}} = \frac{1}{|D_{\text{val}}|}
    \sum_{i=1}^{|D_{\text{val}}|} t_{\text{DL}}(f_i),
\end{equation}
\begin{equation}
    T_{\text{skip}}(\theta) = \frac{1}{|D_{\text{val}}|}
    \sum_{i=1}^{|D_{\text{val}}|} t_{\text{skip}}(f_i, \theta),
\end{equation}
where $f_i$ denotes the $i$-th frame in $D_{\text{val}}$, and both averages are
taken over all frames to reflect the runtime cost on a typical video stream.

% \textbf{Feasibility constraints.} A candidate $\theta$ is considered feasible
% only if it satisfies two hard constraints:
% \begin{equation}
%     R(\theta) \geq R_{\text{base}} - \delta_R,
% \end{equation}
% \begin{equation}
%     T_{\text{skip}}(\theta) \leq \beta \cdot T_{\text{DL}},
% \end{equation}
% where recall tolerance $\delta_R > 0$ is the maximum allowable absolute recall
% drop, and overhead ratio $\beta \in (0, 1)$ bounds the skip module overhead as
% a fraction of one full DL inference. For example, setting $\delta_R = 0.01$ and
% $\beta = 0.10$ means the system tolerates at most a 1\% recall drop and requires
% the skip module to run within 10\% of the cost of a single DL inference.


\textbf{Feasibility constraints.} A candidate $\theta$ is considered feasible
only if it satisfies a hard recall constraint:
\begin{equation}
    R(\theta) \geq R_{\text{base}} - \delta_R,
\end{equation}
where recall tolerance $\delta_R > 0$ is the maximum allowable absolute recall
drop. For example, setting $\delta_R = 0.01$ means the system tolerates at most
a 1\% recall drop compared to the baseline (i.e, pipeline without skip module). 

\textbf{Scoring feasible candidates.}
Among feasible candidates, both objectives are normalized over
$\Theta_{\text{feasible}}$ before weighting, so that the weights
$w_S$ and $w_R$ reflect the intended relative importance rather
than the raw scale of each term.
The range-normalized recall retention and skip rate are defined as:
\begin{equation}
    \tilde{\rho}(\theta)
    = \frac{\rho(\theta) - \min_{\Theta_f}\rho}
           {\max_{\Theta_f}\rho - \min_{\Theta_f}\rho},
    \qquad
    \tilde{S}_r(\theta)
    = \frac{S_r(\theta) - \min_{\Theta_f}S_r}
           {\max_{\Theta_f}S_r  - \min_{\Theta_f}S_r},
\end{equation}
where $\Theta_f \equiv \Theta_{\text{feasible}}$,
and $\rho(\theta)$ is defined as in Equation~\eqref{eq:rho}.
Both normalized terms are bounded in $[0, 1]$ over the feasible
region, ensuring that the weights $w_S$ and $w_R$ directly control
the relative influence of each objective on the final ranking.

\textbf{Optimal selection.}
The optimal parameter set $\theta^*$ is selected as
\begin{equation}
    \theta^*
    = \arg\max_{\theta \in \Theta_{\text{feasible}}}
      \Bigl[
          w_R\,\tilde{\rho}(\theta)
          +
          w_S\,\tilde{S}_r(\theta)
      \Bigr],
\end{equation}
where $\Theta_{\text{feasible}} = \{\theta \in \Theta :
R(\theta) \geq R_{\text{base}} - \delta_R
\ \text{and}\ T_{\text{skip}}(\theta) \leq \beta \cdot T_{\text{DL}}\}$,
and the nonnegative weights satisfy $w_S + w_R = 1$.

\textbf{FAR of the skip-enabled system.} The skip module $\mathcal{S}$ acts as
a conservative gate: skipped frames are labeled negative directly, while passed
frames are processed by the same downstream detector $\mathcal{M}$ as in the
baseline. Therefore, every false alarm produced by the skip-enabled system is
also present in the baseline, implying
\begin{equation}
    \mathrm{FAR}(\theta) \leq \mathrm{FAR}_{\text{base}},
\end{equation}
where $\mathrm{FAR}(\theta)$ and $\mathrm{FAR}_{\text{base}}$ denote the false
alarm rates of the skip-enabled and baseline systems, respectively. The primary
safety risk is recall degradation caused by incorrectly skipped fire/smoke
frames, which motivates enforcing the recall constraint as a hard gate prior to
any candidate ranking.

% !=======================================================
% ! Start justification here 0
% !=======================================================
% ! Start justification here 0
\textbf{Exclusion of FAR-based terms from the scoring objective.}
We derive an assumption-free interval for $\mathrm{FAR}(\theta)$ in terms of
$S_r(\theta)$ and the fixed baseline constant
$\mathrm{FAR}_{\text{base}}$, and then argue that explicitly optimizing
for FAR reduction is unnecessary in the present setting.

By definition:
\begin{equation}\label{eq:far_def}
    \mathrm{FAR}(\theta) = \frac{FP(\theta)}{N_{\text{neg}}}
\end{equation}
where $N_{\text{neg}}$ is the total number of ground-truth negative frames
and $FP(\theta)$ is the number of false positives produced.
Since every skipped frame receives a negative label, false positives arise
only from frames passed to $\mathcal{M}$:
\begin{equation}\label{eq:fp_decomp}
    FP(\theta) = FP_{\text{base}} - a, \quad a \geq 0
\end{equation}
where $a$ is the number of skipped frames that $\mathcal{M}$ would have
misclassified as positive.
A skipped false positive must simultaneously be a baseline false positive
and a skipped negative frame, so:
\begin{equation}\label{eq:a_bound}
    0 \;\leq\; a \;\leq\; \min\!\bigl(FP_{\text{base}},\;N_{\text{skip\_neg}}\bigr)
\end{equation}
This bound requires no assumption on the relationship between the skip
decisions of $\mathcal{S}$ and the outputs of $\mathcal{M}$.
Substituting \eqref{eq:a_bound} into \eqref{eq:fp_decomp}, dividing by
$N_{\text{neg}}$, and using
$\mathrm{FAR}_{\text{base}} = FP_{\text{base}}/N_{\text{neg}}$ and
$S_r(\theta) = N_{\text{skip\_neg}}/N_{\text{neg}}$ gives:
\begin{equation}\label{eq:far_interval}
    \mathrm{FAR}(\theta)
    \;\in\;
    \Bigl[
        \max\!\bigl(0,\;\mathrm{FAR}_{\text{base}} - S_r(\theta)\bigr),
        \;\;
        \mathrm{FAR}_{\text{base}}
    \Bigr]
\end{equation}

\textbf{Justification for excluding FAR optimization.}
Equation~\eqref{eq:far_interval} does not imply that $\mathrm{FAR}(\theta)$
is uniquely determined by $S_r(\theta)$: two configurations with the same
skip rate may skip different frames and therefore yield different realized
values of $a$, leading to different FAR values within the same interval.
However, explicit FAR optimization is not necessary here for two reasons.

First, \eqref{eq:far_interval} provides a hard guarantee that
$\mathrm{FAR}(\theta) \leq \mathrm{FAR}_{\text{base}}$ for every $\theta$.
Thus, the skip module cannot increase the false alarm rate relative to the
baseline.

Second, $\mathcal{S}$ makes skip decisions using only inter-frame motion,
without access to the outputs or error behavior of $\mathcal{M}$.
Therefore, the hyperparameters of $\mathcal{S}$ do not directly control
which baseline false positives are removed.
Any observed FAR reduction depends on whether the motion patterns selected
by $\mathcal{S}$ happen to overlap with frames that $\mathcal{M}$ would
misclassify on the validation set.
Including FAR in the objective could therefore favor configurations that
match validation-specific error patterns, which may not generalize well.

For these reasons, we exclude FAR-based terms from the scoring objective
and report $\mathrm{FAR}(\theta)$ only as an observed consequence metric
in the results tables.